%% Section 6 — Illustrative Case
\section{An Illustrative Comparison: Northern Virginia and Amsterdam}
\label{sec:case}

The following comparison is explicitly illustrative rather than causal. \Cref{tab:comparison} summarises the key dimensions across the two cases; the sections below develop each in turn. Its function is to demonstrate the analytical purchase of the CAS sensemaking framework: to show that the conceptual, structural, and temporal dimensions, examined together and across data modalities, generate an explanation for differential governance outcomes that existing single-method accounts cannot produce. The comparison draws on published sources, secondary reporting, and the preliminary outputs of the WP1--WP4 pipelines run against the project's seed facility index, and it acknowledges this as a limitation to be addressed as the research programme develops.

\subsection{Northern Virginia: The Governance Gap}

Loudoun County, Virginia — marketed, with some accuracy, as ``Data Center Alley'' — is the single largest concentration of data center capacity on earth. Hyperscale campuses operated by Amazon Web Services, Microsoft, Google, and Meta, combined with colocation facilities run by Equinix, Digital Realty, and CyrusOne, have made the county responsible for an estimated 70 per cent of the world's internet traffic at peak. The physical scale is extraordinary: hundreds of server halls, cooling towers visible from the Dulles Toll Road, diesel generator farms staged at the perimeter of campuses whose footprint exceeds many industrial estates.

The \textit{conceptual dimension} of the Northern Virginia assemblage is dominated by the cloud imaginary in its most refined form. Corporate websites for the major operators present the region in terms of network density, carrier diversity, hyperscale capacity, and sustainability commitments that foreground renewable energy sourcing. What is conspicuously absent from this textual register is any acknowledgement of the material friction that has begun to surface in planning hearings, water studies, and local electoral politics: communities in the western parts of the county objecting to the pace of rezoning, state legislators raising questions about water consumption from the Potomac River basin, and municipal planners struggling to translate the county's data center tax revenues into infrastructure capable of absorbing the load. The WP1 divergence analysis flags a high sustainability-to-community ratio and a high ECI gap relative to the Wikipedia register for Ashburn — the characteristic signature of territorial imaginary production operating under low disclosure pressure.

The \textit{structural dimension} is defined by REIT ownership. The majority of leasable data center capacity in Loudoun County is held by publicly traded REITs whose beneficial owners are pension funds, retail investors, and sovereign wealth funds distributed across multiple continents. The WP3 ownership graph for the region reveals a structure in which no single actor — not the operators, not the asset owners, not the investment managers — bears the concentrated accountability that governance intervention requires. The county's planning authority has legal powers over development approvals but no mechanisms to address the governance consequences of dispersed beneficial ownership. The result is a political economy in which the community contestation signal is real, persistent, and growing, but the institutional architecture for translating that signal into regime response is structurally absent. The city2graph urban embeddedness layer reinforces this reading at the spatial scale: the Ashburn facility records 36 power-infrastructure nodes within its 2\,km bounding box but only one residential zone and zero residential-adjacency edges within 1\,km --- a spatial segregation that insulates the assemblage from community entanglement at the built-environment level, not merely at the ownership level.

\subsection{Amsterdam: The Moratorium Episode}

Amsterdam's data center market achieved a global concentration second only to Northern Virginia before the municipality declared a moratorium on new hyperscale development in 2019, effective initially as a pause and then extended through 2022 before being selectively lifted under conditions. The moratorium is a documented governance episode: its triggers, its negotiation, and its resolution are on the record in planning documents, municipal council proceedings, and the ethnographic account of \citet{vonderau2019storing}. It constitutes, for the purposes of this paper, the kind of full transition episode — with identifiable signal, regime response, and outcome — that the CAS temporal dimension is designed to analyse.

The \textit{conceptual dimension} of the Amsterdam assemblage differs from Northern Virginia in ways that are legible in the textual and visual registers. Operators communicating about Amsterdam facilities cannot produce the cloud imaginary without engaging its territorial consequences: the Dutch public and political discourse around data center energy and water consumption, the AMS-IX cooperative that has made the city's internet exchange role a matter of civic identity, and the presence of academic and journalistic attention that has made the infrastructure visible in ways unusual for this sector. The WP1 ECI divergence for Amsterdam is lower than for Loudoun County — though, crucially, whether this reflects genuine territorial embeddedness or a more sophisticated rhetorical absorption of regulatory critique remains an open analytical question, one that the divergence scoring alone cannot resolve but that the close reading the keyword matrix enables can begin to address.

The \textit{structural dimension} is the decisive variable. Amsterdam's data center market in 2019 included a sufficient proportion of operator-owned and cooperative infrastructure — most critically, the AMS-IX member-owned exchange — to create identifiable, accountable counterparties for the municipality's governance intervention. The city's planning system had the legal authority to act; the ownership structure provided entities capable of entering into binding commitments; and the spatial embeddedness of facilities in a dense urban fabric meant that the political constituency for action was geographically concentrated and electorally significant. These are the structural conditions that made the moratorium achievable --- not the intensity of the contestation signal, which was comparable to Northern Virginia, but the architecture of the system receiving it. The urban embeddedness data corroborates: 61 residential zones within 2\,km and 16 within 1\,km proximity means that any governance intervention necessarily negotiates with a politically concentrated and spatially proximate constituency.

\subsection{What the Comparison Demonstrates}

The Virginia--Amsterdam comparison demonstrates, in miniature, the analytical programme of the paper. The same external pressure — community and political contestation of data center development — produces structurally different outcomes across two contexts. The variable that differs is not signal strength but the CAS properties of the receiving system: the ownership architecture (concentrated and accountable in Amsterdam, dispersed and insulated in Virginia), the planning authority's governance capacity (legally empowered and spatially concentrated in Amsterdam, constrained and revenue-dependent in Virginia), and the territorial imaginary (contested and partially negotiated in Amsterdam, dominant and entrenched in Virginia). These are the structural and conceptual dimensions of the assemblage, and they are only visible simultaneously — the ownership graph alone cannot explain why the municipality could act, the text analysis alone cannot explain why the ownership structure mattered.

This is the argument for multimodality: not that more data is better, but that different data modalities reach different dimensions of the assemblage, and that no single dimension is explanatorily self-sufficient.
